\chapter{پاسخ تمرین‌های منتخب}
\label{app:solutions}

در این پیوست پاسخ شماری از تمرین‌های کتاب آمده است. پاسخ‌ها یک راه حل‌اند،
نه تنها راه حل؛ اگر برنامه‌ی شما خروجی درست می‌دهد و خوانا است، درست است،
حتی اگر با این‌جا فرق داشته باشد. خواهش می‌کنیم پیش از نگاه کردن به هر
پاسخ، خودتان تلاش کنید.

\section*{فصل \ref{ch:first}}

\subsection*{تمرین ۱.۱}
\begin{salam}
تابع اصلی:
    چاپ "سارا"
    چاپ "احمدی"
    چاپ "کاشان"
تمام
\end{salam}

\subsection*{تمرین ۱.۳}
\begin{salam}
تابع اصلی:
    چاپ "حاصل جمع ۱۲۵ و ۳۷۵:", 125 + 375
    چاپ "حاصل ضرب ۲۴ در ۷:", 24 * 7
    چاپ "ساعت‌های یک هفته:", 7 * 24
تمام
\end{salam}
\begin{outputfa}
حاصل جمع ۱۲۵ و ۳۷۵: 500
حاصل ضرب ۲۴ در ۷: 168
ساعت‌های یک هفته: 168
\end{outputfa}

\subsection*{تمرین ۱.۴}
سه اشتباه: دونقطه‌ی بعد از \fcode{تابع اصلی} فراموش شده، در خط دوم به
جای \code{"} از گیومه‌ی فارسی استفاده شده، و رشته‌ی خط سوم نقل قول پایانی
ندارد.

\section*{فصل \ref{ch:variables}}

\subsection*{تمرین ۲.۱}
\begin{salam}
تابع اصلی:
    طول := 5
    عرض := 4
    ارتفاع := 3
    چاپ "مساحت کف:", طول * عرض, "متر مربع"
    چاپ "حجم اتاق:", طول * عرض * ارتفاع, "متر مکعب"
تمام
\end{salam}
\begin{outputfa}
مساحت کف: 20 متر مربع
حجم اتاق: 60 متر مکعب
\end{outputfa}

\subsection*{تمرین ۲.۳}
خروجی \code{7 1} است. دنبال کنید: \fcode{الف} برابر ۳ و \fcode{ب} برابر
۶ می‌شود. سپس \fcode{الف} برابر ۷ می‌شود. در آخر \fcode{ب} برابر
\code{7 - 6} یعنی ۱ می‌شود.

\subsection*{تمرین ۲.۴}
\fcode{امتیاز} بدون \fcode{متغیر} تعریف شده اما در خط بعد عوض می‌شود؛ باید
\fcode{متغیر امتیاز \lr{:=} 10} باشد. \fcode{نام بازیکن} تعریف شده اما
استفاده نشده؛ یا در \fcode{چاپ} به کارش ببرید یا خطش را حذف کنید.

\section*{فصل \ref{ch:types}}

\subsection*{تمرین ۳.۳}
\begin{salam}
تابع اصلی:
    دقیقه ها := 135
    ساعت := دقیقه ها / 60
    دقیقه := دقیقه ها % 60
    چاپ ساعت, "ساعت و", دقیقه, "دقیقه"
تمام
\end{salam}
\begin{outputfa}
2 ساعت و 15 دقیقه
\end{outputfa}

\section*{فصل \ref{ch:operators}}

\subsection*{تمرین ۴.۲}
\begin{salam}
تابع اصلی:
    عدد := 4729
    یکان := عدد % 10
    دهگان := عدد / 10 % 10
    صدگان := عدد / 100 % 10
    هزارگان := عدد / 1000
    چاپ "مجموع رقم‌ها:", یکان + دهگان + صدگان + هزارگان
تمام
\end{salam}
\begin{outputfa}
مجموع رقم‌ها: 22
\end{outputfa}

\section*{فصل \ref{ch:conditions}}

\subsection*{تمرین ۵.۳}
\begin{salam}
تابع اصلی:
    ساعت := 15
    اگر ساعت < 12:
        چاپ "صبح بخیر"
    وگرنه ساعت < 17:
        چاپ "ظهر بخیر"
    وگرنه ساعت < 20:
        چاپ "عصر بخیر"
    وگرنه:
        چاپ "شب بخیر"
    تمام
تمام
\end{salam}
\begin{outputfa}
ظهر بخیر
\end{outputfa}

\subsection*{تمرین ۵.۶}
سه خط چاپ می‌شود: «بزرگ‌تر از ده»، «مضرب پنج» و «مضرب سه». زنجیره‌ی اول
با نخستین شرط درست تمام می‌شود و «بزرگ‌تر از پنج» چاپ نمی‌شود؛ اما دو
\fcode{اگر} بعدی جداگانه‌اند و هر دو بررسی می‌شوند.

\section*{فصل \ref{ch:loops}}

\subsection*{تمرین ۶.۱}
\begin{salam}
تابع اصلی:
    متغیر خط := ""
    تکرار 1 تا 50 با عدد:
        اگر عدد % 7 == 0:
            خط = خط + عدد + " "
        تمام
    تمام
    چاپ خط
تمام
\end{salam}
\begin{outputfa}
7 14 21 28 35 42 49
\end{outputfa}

\subsection*{تمرین ۶.۲}
\begin{salam}
تابع اصلی:
    متغیر حاصل := 1
    تکرار 1 تا 10 با عدد:
        حاصل *= عدد
    تمام
    چاپ حاصل
تمام
\end{salam}
\begin{outputfa}
3628800
\end{outputfa}

\subsection*{تمرین ۶.۶}
شش خط. حلقه‌ی بیرونی سه دور دارد و در هر دور، حلقه‌ی درونی فقط برای
\fcode{ب} برابر ۰ و ۱ چاپ می‌کند و در ۲ می‌شکند.

\section*{فصل \ref{ch:functions}}

\subsection*{تمرین ۷.۱}
\begin{salam}
تابع بزرگتر(الف: صحیح, ب: صحیح): صحیح:
    اگر الف > ب: بازگشت الف تمام
    بازگشت ب
تمام

تابع اصلی:
    چاپ بزرگتر(بزرگتر(12, 25), 19)
تمام
\end{salam}
\begin{outputfa}
25
\end{outputfa}
بزرگ‌ترین سه عدد، بزرگ‌تر از «بزرگ‌تر دو تای اول» و سومی است.

\subsection*{تمرین ۷.۶}
وقتی \fcode{عدد} صفر است، هیچ‌کدام از دو شرط برقرار نیست و تابع بدون
\fcode{بازگشت} تمام می‌شود؛ سلام این را خطا می‌گیرد. راه حل: شرط دوم را
به \fcode{وگرنه:} بی‌شرط تبدیل کنید، یا در پایان تابع
\fcode{بازگشت عدد} بگذارید.

\section*{فصل \ref{ch:arrays}}

\subsection*{تمرین ۸.۲}
\begin{salam}
تابع میانگین(اعداد: صحیح[:]): اعشار۶۴:
    متغیر جمع := 0
    هر عدد در اعداد:
        جمع += عدد
    تمام
    تعداد := اعداد.طول()
    بازگشت (جمع بعنوان اعشار۶۴) / (تعداد بعنوان اعشار۶۴)
تمام

تابع بیشترین(اعداد: صحیح[:]): صحیح:
    متغیر بزرگترین := اعداد[0]
    هر عدد در اعداد:
        اگر عدد > بزرگترین: بزرگترین = عدد تمام
    تمام
    بازگشت بزرگترین
تمام

تابع اصلی:
    نمرات := [17, 15, 19, 12, 20]
    چاپ "میانگین:", میانگین(نمرات[:])
    چاپ "بیشترین:", بیشترین(نمرات[:])
تمام
\end{salam}
\begin{outputfa}
میانگین: 16.6
بیشترین: 20
\end{outputfa}

\subsection*{تمرین ۸.۶}
خروجی \code{[10, 2, 30, 4, 50]} است: عضوهایی که اندیس زوج دارند (۰، ۲
و ۴) ده برابر می‌شوند.

\section*{فصل \ref{ch:strings}}

\subsection*{تمرین ۹.۴}
\begin{salam}
تابع اصلی:
    متن := "17,15,19,12"
    تکه ها := متن.بشکاف(",")
    متغیر مجموع := 0
    تکرار تکه ها.طول() با ک:
        مجموع += تکه ها.بگیر(ک).به صحیح()
    تمام
    تعداد := تکه ها.طول()
    چاپ "مجموع:", مجموع
    چاپ "میانگین:", (مجموع بعنوان اعشار۶۴) / (تعداد بعنوان اعشار۶۴)
تمام
\end{salam}
\begin{outputfa}
مجموع: 63
میانگین: 15.75
\end{outputfa}
